\chapter{Worked Problem Laboratories}
\label{chap:worked-problem-labs}

\section{Purpose And Scope}
\label{sec:worked-problem-labs-purpose}

This chapter collects self-contained lecture-note laboratories. Each laboratory
has a precise model, a defined notation set, and a derivation that can be read
from the definitions and assumptions given here. The examples are intentionally
varied: they show how the abstract machinery in the earlier chapters behaves
when used on concrete mechanics problems.

\begin{table}[htbp]
\centering
\caption{Notation used in the worked problem laboratories. Local definitions in
each section override this table when stated.}
\label{tab:worked-lab-notation}
\small
\begin{tabular}{p{0.18\textwidth}p{0.48\textwidth}p{0.25\textwidth}}
\toprule
Symbol & Meaning & Type\\
\midrule
\(s^a\) & new generalized coordinates under a point transformation & coordinates on \(Q\)\\
\(q^i\) & old generalized coordinates & coordinates on \(Q\)\\
\(V(x)\) & one-dimensional potential energy & scalar function\\
\(\gamma\) & damping coefficient in the oscillator example & inverse time\\
\(\omega\) & imposed angular speed in rotating-hoop example & inverse time\\
\(a\) & hoop radius or semimajor axis depending on section & locally defined\\
\(R,r\) & cylinder radii or Schwarzschild length/radius depending on section & locally defined\\
\(N\) & normal force in rolling-contact examples & force\\
\(A(r)\) & \(1-R/r\) in the Schwarzschild-type example & dimensionless\\
\(E,L_z\) & conserved energy and angular momentum in central examples & constants of motion\\
\(H_{\mathrm{HH}}\) & Henon-Heiles Hamiltonian & Hamiltonian function\\
\(\lambda_{\mathrm{HH}}\) & Henon-Heiles perturbation strength & small parameter\\
\(I_i,\theta_i\) & oscillator actions and angles & canonical coordinates\\
\(r_i,v_i\) & positions and velocities of three equal masses & vectors in \(\R^2\)\\
\(m,k\) & mass and Newtonian coupling in three-body lab & positive constants\\
\(\alpha,\beta\) & internal rotor angles in the helicopter lab & shape coordinates\\
\(\mathcal A_\alpha,\mathcal A_\beta\) & gauge connection components over rotor shape space & \(\mathfrak{so}(3)\)-valued\\
\bottomrule
\end{tabular}
\end{table}

\paragraph{How to use this chapter.}
Each laboratory is deliberately small, but each one isolates a recurring
mechanics move: change coordinates, test regularity, encode dissipation with
time dependence, reduce a constrained rolling problem, test circular-orbit
stability, expose resonance, or reconstruct motion from a gauge connection.
The value of the chapter is cumulative. A reader should be able to take the
workflow of one section and recognize the same workflow in a more complicated
system later.

Each laboratory should be read with three questions in mind. What is the
minimal model? What calculation turns the model into a diagnostic quantity?
What would fail if one of the assumptions were changed? These questions are
what turn a worked example into a reusable method rather than a solved problem
to memorize.

\paragraph{How figures function in this chapter.}
The figures are not decorative summaries after the calculation. Each one is
attached to a modelling decision: a coordinate choice, a constraint force, an
effective potential, a phase-space section, an initial-data geometry, or a
gauge loop. When using these laboratories as exercise problems, the figure
should be read before the algebra and then checked again after the final boxed
result. If a student cannot say which equation the figure represents, the
problem has not yet been understood geometrically.

\section{Form-Invariance Of Lagrange Equations}
\label{sec:lab-point-transformations}

We begin with the coordinate-independence of the Euler-Lagrange equations. Let
\[
  q^i=q^i(s^1,\ldots,s^n,t)
\]
be a regular time-dependent point transformation. Regular means that the
Jacobian \(\partial q^i/\partial s^a\) has rank \(n\). The transformed
Lagrangian is
\[
  L'(s,\dot s,t)=L(q(s,t),\dot q(s,\dot s,t),t),
\]
where
\[
  \dot q^i
  =
  \frac{\partial q^i}{\partial s^a}\dot s^a
  +
  \frac{\partial q^i}{\partial t}.
\]
Two elementary identities do the work:
\[
  \frac{\partial \dot q^i}{\partial \dot s^a}
  =
  \frac{\partial q^i}{\partial s^a},
  \qquad
  \frac{\dd}{\dd t}
  \left(
    \frac{\partial q^i}{\partial s^a}
  \right)
  =
  \frac{\partial \dot q^i}{\partial s^a}.
\]
Using the chain rule,
\[
  \frac{\partial L'}{\partial \dot s^a}
  =
  \frac{\partial L}{\partial \dot q^i}
  \frac{\partial q^i}{\partial s^a},
\]
and
\[
  \frac{\partial L'}{\partial s^a}
  =
  \frac{\partial L}{\partial q^i}
  \frac{\partial q^i}{\partial s^a}
  +
  \frac{\partial L}{\partial \dot q^i}
  \frac{\partial \dot q^i}{\partial s^a}.
\]
Therefore
\[
\begin{aligned}
  \frac{\dd}{\dd t}\frac{\partial L'}{\partial \dot s^a}
  -
  \frac{\partial L'}{\partial s^a}
  &=
  \left(
    \frac{\dd}{\dd t}\frac{\partial L}{\partial\dot q^i}
    -
    \frac{\partial L}{\partial q^i}
  \right)
  \frac{\partial q^i}{\partial s^a}.
\end{aligned}
\]
If \(q(t)\) satisfies the Euler-Lagrange equations, the right side vanishes;
if the transformation is regular, the converse holds as well. Thus the
Euler-Lagrange equation is not tied to a preferred coordinate system.

\section{A Nonstandard Lagrangian And Regularity}
\label{sec:lab-nonstandard-lagrangian}

This laboratory separates two ideas that are often conflated. A Lagrangian may
produce a familiar-looking equation of motion and still fail to be globally
regular. Regularity is a statement about whether velocities can be solved from
momenta. Without it, the Legendre transform can fold phase space and the
Hamiltonian description must be handled with care.

Consider the one-dimensional Lagrangian
\[
  L(x,\dot x)
  =
  \frac{m^2\dot x^4}{12}
  +
  m\dot x^2V(x)
  -
  V(x)^2.
\]
The momentum is
\[
  p=\frac{\partial L}{\partial\dot x}
  =
  \frac{m^2}{3}\dot x^3+2mV\dot x.
\]
The Euler-Lagrange equation gives
\[
\begin{aligned}
  0
  &=
  \frac{\dd}{\dd t}
  \left(
    \frac{m^2}{3}\dot x^3+2mV\dot x
  \right)
  -
  \left(
    m\dot x^2V'-2VV'
  \right)\\
  &=
  (m\dot x^2+2V)(m\ddot x+V').
\end{aligned}
\]
There are two branches. The Newtonian branch is
\[
  m\ddot x=-V'(x).
\]
The second branch,
\[
  m\dot x^2+2V(x)=0,
\]
is a singular energy-surface condition. The velocity Hessian
\[
  \frac{\partial^2L}{\partial\dot x^2}
  =
  m(m\dot x^2+2V)
\]
vanishes on this second branch. This example is a useful warning: a Lagrangian
can produce familiar equations on a regular branch while being singular on
another branch. Regularity is not decorative; it is what lets the Legendre
transform define a Hamiltonian system.

\begin{figure}[htbp]
\centering
\begin{tikzpicture}[scale=1.0, >=Latex]
  \draw[->] (-3.0,0) -- (3.15,0) node[right] {\(\dot x\)};
  \draw[->] (0,-2.2) -- (0,2.25) node[above] {\(p\)};
  \draw[blue!70!black, thick, domain=-2.55:2.55, samples=180]
    plot (\x,{0.33*\x*\x*\x-0.9*\x});
  \draw[dashed, red!75!black] (-0.95,-1.9) -- (-0.95,1.9);
  \draw[dashed, red!75!black] (0.95,-1.9) -- (0.95,1.9);
  \fill[red!75!black] (-0.95,0.57) circle (1.3pt)
    node[above left] {\(W=0\)};
  \fill[red!75!black] (0.95,-0.57) circle (1.3pt)
    node[below right] {\(W=0\)};
  \draw[decorate, decoration={brace, amplitude=5pt}]
    (-0.95,-1.15) -- (0.95,-1.15)
    node[midway, below=7pt] {multi-valued velocity branch};
  \node[blue!70!black, align=center] at (2.05,1.45)
    {folded\\Legendre map};
\end{tikzpicture}
\caption{Nonstandard Lagrangian laboratory. For a representative region with
\(V<0\), the momentum-velocity map can fold where the velocity Hessian
\(W=\partial^2L/\partial\dot x^2\) vanishes. At such points the Legendre
transform is not locally invertible, even though the Euler-Lagrange equation has
a familiar Newtonian branch.}
\label{fig:lab-nonstandard-legendre-fold}
\end{figure}
\FloatBarrier

\section{Damped Oscillator From A Time-Dependent Lagrangian}
\label{sec:lab-damped-oscillator}

Let
\[
  L(q,\dot q,t)
  =
  e^{\gamma t}
  \left(
    \frac12m\dot q^2-\frac12kq^2
  \right).
\]
The Euler-Lagrange equation is
\[
  \frac{\dd}{\dd t}(e^{\gamma t}m\dot q)+e^{\gamma t}kq=0,
\]
or
\[
  \boxed{
  m\ddot q+m\gamma\dot q+kq=0.
  }
\]
Although this is a dissipative equation, it comes from a time-dependent
Lagrangian. The canonical energy is not conserved because \(L\) depends
explicitly on time.

Set
\[
  s=e^{\gamma t/2}q,
  \qquad
  q=e^{-\gamma t/2}s.
\]
Then
\[
  \dot q=e^{-\gamma t/2}\left(\dot s-\frac{\gamma}{2}s\right).
\]
Substituting into the action gives
\[
  L
  =
  \frac12m\left(\dot s-\frac{\gamma}{2}s\right)^2
  -
  \frac12ks^2.
\]
The cross term is a total derivative:
\[
  -\frac12m\gamma s\dot s
  =
  -\frac{\dd}{\dd t}\left(\frac14m\gamma s^2\right).
\]
Discarding that total derivative, the equivalent Lagrangian is
\[
  L_{\mathrm{eff}}
  =
  \frac12m\dot s^2
  -
  \frac12m
  \left(
    \frac{k}{m}-\frac{\gamma^2}{4}
  \right)s^2.
\]
Thus
\[
  \ddot s+\left(\frac{k}{m}-\frac{\gamma^2}{4}\right)s=0.
\]
The transformation has moved damping into a shifted oscillator frequency.

\begin{figure}[htbp]
\centering
\begin{tikzpicture}[scale=1.0, >=Latex]
  \draw[->] (-0.2,0) -- (6.45,0) node[right] {\(t\)};
  \draw[->] (0,-1.45) -- (0,1.6) node[above] {amplitude};
  \draw[blue!70!black, dashed, domain=0:6.0, samples=120]
    plot (\x,{1.10*exp(-0.22*\x)});
  \draw[blue!70!black, dashed, domain=0:6.0, samples=120]
    plot (\x,{-1.10*exp(-0.22*\x)});
  \draw[blue!70!black, thick, domain=0:6.0, samples=180]
    plot (\x,{1.10*exp(-0.22*\x)*cos(2.7*\x r)});
  \draw[red!75!black, thick, domain=0:6.0, samples=180]
    plot (\x,{1.10*cos(2.7*\x r)});
  \node[blue!70!black] at (4.7,0.43) {\(q(t)\)};
  \node[red!75!black] at (4.5,1.25) {\(s=e^{\gamma t/2}q\)};
  \node[align=center] at (3.2,-1.85)
    {rescaling absorbs the exponential damping envelope};
\end{tikzpicture}
\caption{Damped-oscillator laboratory. In the original coordinate \(q\), the
motion has a decaying envelope. The rescaled coordinate \(s=e^{\gamma t/2}q\)
obeys an undamped oscillator equation with shifted frequency.}
\label{fig:lab-damped-oscillator-rescaling}
\end{figure}
\FloatBarrier

\section{Rotating Hoop: Effective Potential And Bifurcation}
\label{sec:lab-rotating-hoop}

A bead of mass \(m\) moves without friction on a hoop of radius \(a\). The hoop
is vertical and rotates about its vertical diameter with imposed angular speed
\(\omega\). Let \(\theta\) be the angle of the bead measured from the downward
vertical. The bead position has speed
\[
  v^2=a^2\dot\theta^2+a^2\omega^2\sin^2\theta.
\]
The Lagrangian is
\[
  L
  =
  \frac12ma^2\dot\theta^2
  +
  \frac12ma^2\omega^2\sin^2\theta
  +
  mga\cos\theta.
\]

\begin{figure}[htbp]
\centering
\begin{tikzpicture}[scale=1.0, >=Latex]
  \begin{scope}
    \draw[->] (0,-2.05) -- (0,2.25) node[above] {rotation axis};
    \draw[dashed] (0,-1.85) -- (0,1.85);
    \draw[thick] (0,0) circle (1.65);
    \draw[fill=red!70!black] (1.17,-1.17) circle (0.09);
    \draw[thick] (0,0) -- (1.17,-1.17);
    \draw[->] (0,-0.75) arc (-90:-45:0.75);
    \node at (0.55,-0.95) {\(\theta\)};
    \draw[->, blue!70!black] (1.17,-1.17) -- (1.75,-1.17)
      node[right] {azimuthal speed \(a\omega\sin\theta\)};
    \node[below] at (0,-1.72) {downward vertical};
    \node at (0,-2.45) {bead on rotating vertical hoop};
  \end{scope}
  \begin{scope}[xshift=5.0cm]
    \draw[->] (-2.0,0) -- (2.0,0) node[right] {\(\theta\)};
    \draw[->] (0,-1.5) -- (0,1.55) node[above] {\(V_{\mathrm{eff}}\)};
    \draw[gray!60] (-1.55,-0.05) -- (-1.55,0.05) node[below=3pt] {\(-\pi\)};
    \draw[gray!60] (1.55,-0.05) -- (1.55,0.05) node[below=3pt] {\(\pi\)};
    \draw[blue!70!black, thick, domain=-1.65:1.65, samples=140]
      plot (\x,{-0.35*sin(2*\x r)*sin(2*\x r)-0.80*cos(2*\x r)});
    \draw[red!75!black, thick, domain=-1.65:1.65, samples=140]
      plot (\x,{-0.95*sin(2*\x r)*sin(2*\x r)-0.45*cos(2*\x r)+0.25});
    \node[blue!70!black, align=center] at (-1.25,1.15)
      {\(\omega<\omega_c\)\\one minimum};
    \node[red!75!black, align=center] at (1.2,1.15)
      {\(\omega>\omega_c\)\\two minima};
  \end{scope}
\end{tikzpicture}
\caption{Rotating-hoop laboratory. The geometry fixes the kinetic term
\(a^2\omega^2\sin^2\theta\), and the effective potential changes from one
bottom minimum to two off-bottom minima when \(\omega\) crosses
\(\omega_c=\sqrt{g/a}\). The potential curves are schematic and emphasize the
bifurcation, not numerical scale.}
\label{fig:lab-rotating-hoop}
\end{figure}
\FloatBarrier

This is a one-dimensional system with effective potential
\[
  V_{\mathrm{eff}}(\theta)
  =
  -\frac12ma^2\omega^2\sin^2\theta
  -
  mga\cos\theta.
\]
Equilibria satisfy
\[
  \frac{\dd V_{\mathrm{eff}}}{\dd\theta}
  =
  ma\sin\theta(g-a\omega^2\cos\theta)=0.
\]
Thus either
\[
  \theta=0,\pi,
\]
or
\[
  \cos\theta=\frac{g}{a\omega^2}.
\]
The off-bottom equilibria exist only when
\[
  \omega\geq\omega_c=\sqrt{\frac{g}{a}}.
\]
At \(\omega=\omega_c\), the bottom equilibrium changes stability and two
symmetric equilibria appear. This is a concrete mechanics realization of a
pitchfork bifurcation in an effective potential.

\section{Rolling Contact: Contact Forces As Multipliers}
\label{sec:lab-rolling-contact}

Rolling examples force one to distinguish coordinates, constraints, and
reaction forces.

\subsection{Hoop Rolling On A Fixed Cylinder}

A hoop of mass \(m\) and radius \(r\) rolls without slipping on the outside of a
fixed cylinder of radius \(R\). Let \(\theta\) be the angular position of the
hoop center measured from the top. The center of mass moves on a circle of
radius \(R+r\). The rolling condition gives
\[
  r\dot\phi=(R+r)\dot\theta,
\]
where \(\phi\) is the spin angle of the hoop. Since the hoop has
\[
  I_{\mathrm{hoop}}=mr^2,
\]
the kinetic energy is
\[
  T=
  \frac12m(R+r)^2\dot\theta^2
  +
  \frac12mr^2\dot\phi^2
  =
  m(R+r)^2\dot\theta^2.
\]
Starting from rest at the top, energy gives
\[
  m(R+r)^2\dot\theta^2
  =
  mg(R+r)(1-\cos\theta).
\]
The radial equation is
\[
  mg\cos\theta-N=m(R+r)\dot\theta^2,
\]
where \(N\) is the normal force. Loss of contact occurs at \(N=0\). Using the
energy relation,
\[
  \cos\theta=1-\cos\theta,
\]
so
\[
  \boxed{\cos\theta=\frac12.}
\]
The hoop leaves the cylinder at \(\theta=\pi/3\) from the top.

\begin{figure}[htbp]
\centering
\begin{tikzpicture}[scale=1.0, >=Latex]
  \draw[thick, fill=gray!8] (0,0) circle (1.4);
  \draw[thick, fill=blue!7] (1.65,1.65) circle (0.45);
  \draw[dashed] (0,0) -- (1.65,1.65);
  \draw[dashed] (0,0) -- (0,2.35);
  \draw[->] (0,1.9) arc (90:45:1.9);
  \node at (0.85,1.95) {\(\theta\)};
  \draw[->, red!75!black, thick] (1.65,1.65) -- (2.15,2.15)
    node[right] {\(N\)};
  \draw[->, black!70, thick] (1.65,1.65) -- (1.65,0.75)
    node[below] {\(mg\)};
  \draw[->, blue!70!black, thick] (2.1,1.65) arc (0:65:0.45)
    node[right] {\(\dot\phi\)};
  \node[align=center] at (0,-1.95)
    {loss of contact occurs when the\\normal force reaches \(N=0\)};
  \node at (-1.1,-1.1) {fixed cylinder \(R\)};
  \node at (2.7,2.25) {rolling hoop \(r\)};
\end{tikzpicture}
\caption{Rolling-contact laboratory. The radial force balance must be kept
until the end because the contact force \(N\) is the diagnostic for leaving the
surface. Eliminating the rolling constraint gives the motion, but \(N=0\)
decides the loss-of-contact angle.}
\label{fig:lab-rolling-contact}
\end{figure}
\FloatBarrier

\subsection{Solid Cylinder Inside A Hollow Cylinder}

Now consider a solid cylinder of radius \(R\) rolling inside a fixed hollow
cylinder of radius \(2R\). The center of mass moves on a circle of radius
\(R\). With rolling imposed, the kinetic energy is
\[
  T
  =
  \frac12MR^2\dot\theta^2
  +
  \frac14MR^2\dot\theta^2,
\]
because a solid cylinder has moment of inertia \(I=(1/2)MR^2\). Taking the
potential as \(V=-MgR\cos\theta\), the Lagrangian is
\[
  L=\frac34MR^2\dot\theta^2+MgR\cos\theta.
\]
The equation of motion is
\[
  \frac32MR^2\ddot\theta+MgR\sin\theta=0,
\]
or
\[
  \frac32\ddot\theta+\frac{g}{R}\sin\theta=0.
\]
For small oscillations,
\[
  \boxed{
  \omega_{\mathrm{small}}=\sqrt{\frac{2g}{3R}}.
  }
\]

\section{Schwarzschild-Type Central-Force Stability}
\label{sec:lab-schwarzschild-stability}

Consider the central Lagrangian
\[
  L=-m\sqrt{
    A(r)-\frac{\dot r^2}{A(r)}-r^2\dot\theta^2
  },
  \qquad
  A(r)=1-\frac{R}{r},
  \qquad r>R.
\]
This model is a relativistic-looking laboratory for circular-orbit stability.
Define
\[
  B=A-\frac{\dot r^2}{A}-r^2\dot\theta^2.
\]
The conserved energy and angular momentum are
\[
  E=\frac{mA}{\sqrt B},
  \qquad
  L_z=\frac{mr^2\dot\theta}{\sqrt B}.
\]
Eliminating \(\dot\theta\) and \(B\) gives the radial equation in effective
potential form:
\[
  \dot r^2
  =
  \frac{A^2}{E^2}
  \left[
    E^2-A\left(m^2+\frac{L_z^2}{r^2}\right)
  \right].
\]
Thus circular orbits are critical points of
\[
  U_{\mathrm{eff}}(r)
  =
  A(r)\left(m^2+\frac{L_z^2}{r^2}\right)
\]
at fixed \(L_z\). The circular-orbit condition \(U_{\mathrm{eff}}'(r_0)=0\)
gives
\[
  L_z^2
  =
  \frac{Rm^2r_0^2}{2r_0-3R}.
\]
Hence circular orbits require \(r_0>3R/2\). Stability is determined by
\[
  U_{\mathrm{eff}}''(r_0)
  =
  \frac{2Rm^2(r_0-3R)}
  {r_0^3(2r_0-3R)}.
\]
Therefore
\[
  \boxed{
  r_0>3R \text{ is stable},\qquad
  \frac32R<r_0<3R \text{ is unstable}.
  }
\]
This is the same structural calculation used in more realistic relativistic
orbit problems: circular-orbit stability is a second-derivative question for
an effective potential at fixed angular momentum.

\begin{figure}[htbp]
\centering
\begin{tikzpicture}[scale=1.0, >=Latex]
  \draw[->] (0,0) -- (6.4,0) node[right] {\(r/R\)};
  \draw[->] (0,-0.35) -- (0,3.0) node[above] {\(U_{\mathrm{eff}}\)};
  \draw[dashed] (1.0,-0.2) -- (1.0,2.8) node[above] {\(R\)};
  \draw[dashed] (3.0,-0.2) -- (3.0,2.8) node[above] {\(3R\)};
  \draw[red!75!black, thick, domain=1.18:5.9, samples=180]
    plot (\x,{1.15+1.7*(1-1/\x)*(1+2.9/(\x*\x))-0.33*\x});
  \fill[red!75!black] (2.02,1.83) circle (1.3pt)
    node[above right] {unstable circular orbit};
  \fill[blue!70!black] (4.35,1.54) circle (1.3pt)
    node[below right] {stable circular orbit};
  \draw[blue!70!black, thick, ->] (4.05,1.25) -- (4.30,1.48);
  \draw[red!75!black, thick, ->] (1.8,1.55) -- (2.0,1.78);
  \node[align=center] at (3.25,-0.75)
    {stability is the sign of \(U_{\mathrm{eff}}''(r_0)\)\\
     at fixed angular momentum};
\end{tikzpicture}
\caption{Effective-potential picture for the Schwarzschild-type laboratory.
Circular orbits are critical points; the outer critical point is a minimum
\((r_0>3R)\), while the inner one is a maximum
\((3R/2<r_0<3R)\). The curve is schematic.}
\label{fig:lab-schwarzschild-effective-potential}
\end{figure}
\FloatBarrier

\section{Henon-Heiles As The Finite-Dimensional Chaos Bridge}
\label{sec:lab-henon-heiles}

The Henon-Heiles Hamiltonian is
\[
  H_{\mathrm{HH}}
  =
  H_0+\lambda_{\mathrm{HH}}H_1,
\]
with
\[
  H_0=
  \frac12(p_x^2+p_y^2+x^2+y^2),
  \qquad
  H_1=x^2y-\frac13y^3.
\]
For the isotropic oscillator, introduce action-angle variables
\[
  x=\sqrt{2I_1}\sin\theta_1,
  \qquad
  p_x=\sqrt{2I_1}\cos\theta_1,
\]
\[
  y=\sqrt{2I_2}\sin\theta_2,
  \qquad
  p_y=\sqrt{2I_2}\cos\theta_2.
\]
Then
\[
  H_0=I_1+I_2,
  \qquad
  \omega_0=(1,1).
\]
The first-order homological equation for a generating function \(S_1\) is
\[
  \omega_0\cdot\partial_\theta S_1
  =
  \langle H_1\rangle-H_1.
\]
In Fourier components,
\[
  i(k_1+k_2)S_{1k}=-H_{1k}.
\]
The denominator vanishes whenever
\[
  k_1+k_2=0.
\]
At first order the cubic perturbation has no average and the dangerous modes
can be handled by symmetry, but at second order the canonical transformation
generates resonant Fourier modes with \(k_1+k_2=0\). Those modes cannot be
removed by a periodic \(S_2\). This is the small-divisor obstruction in its
most concrete classroom form: the unperturbed oscillator has equal
frequencies, so resonant combinations of angles do not average away.

The Henon-Heiles system is therefore the right finite-dimensional bridge from
integrability to Hamiltonian chaos. At low energy, Poincare sections show
deformed invariant curves. At higher energy, resonance islands and chaotic
regions appear. Unlike the asteroid problem, all variables and forces are
visible in a four-dimensional phase space.

\begin{figure}[htbp]
\centering
\begin{tikzpicture}[scale=1.0, >=Latex]
  \begin{scope}
    \draw[->] (-2.1,0) -- (2.1,0) node[right] {\(x\)};
    \draw[->] (0,-1.9) -- (0,2.2) node[above] {\(y\)};
    \foreach \s in {0.45,0.75,1.05,1.35}
    {
      \draw[blue!70!black, thick, domain=0:360, samples=120, variable=\t]
        plot ({\s*cos(\t)*(1+0.12*sin(3*\t))},
              {\s*sin(\t)*(1-0.10*cos(3*\t))});
    }
    \draw[red!75!black, thick]
      (0,1.75) -- (-1.55,-0.95) -- (1.55,-0.95) -- cycle;
    \node[red!75!black, align=center] at (0,-1.45)
      {three escape\\valleys};
    \node at (0,2.55) {configuration contours};
  \end{scope}
  \begin{scope}[xshift=5.2cm]
    \draw[->] (-1.8,0) -- (1.9,0) node[right] {\(x\)};
    \draw[->] (0,-1.8) -- (0,1.9) node[above] {\(p_x\)};
    \draw[blue!70!black, thick] (0,0) ellipse (1.2 and 0.65);
    \draw[blue!70!black, thick] (0,0) ellipse (0.72 and 0.38);
    \foreach \a in {0,120,240}
      \draw[red!75!black, thick] ({0.82*cos(\a)},{0.82*sin(\a)})
        ellipse (0.28 and 0.13);
    \foreach \x/\y in {-1.3/0.9,-1.1/-0.8,-0.4/1.15,0.25/-1.2,1.0/0.75,1.25/-0.55}
      \fill[black!55] (\x,\y) circle (0.9pt);
    \node at (0,2.35) {Poincare-section cartoon};
  \end{scope}
\end{tikzpicture}
\caption{Henon-Heiles laboratory. Configuration-space contours show the
threefold cubic deformation of the oscillator potential; a Poincare section
then records the phase-space transition from invariant curves to islands and
chaotic scatter. Both panels are schematic diagnostics for the derivation.}
\label{fig:lab-henon-heiles}
\end{figure}
\FloatBarrier

\section{Equal-Mass Three-Body Laboratory}
\label{sec:lab-three-body}

For an equal-mass three-body laboratory, take three masses in the plane with
pair potential
\[
  V=-k\left(\frac1{r_{12}}+\frac1{r_{23}}+\frac1{r_{13}}\right),
\]
where
\[
  r_{ij}=\norm{r_i-r_j}.
\]
The Lagrangian is
\[
  L=\sum_{i=1}^3\frac12m\norm{\dot r_i}^2
  +
  k\left(\frac1{r_{12}}+\frac1{r_{23}}+\frac1{r_{13}}\right).
\]
The equations are
\[
  m\ddot r_i
  =
  -k\sum_{j\neq i}\frac{r_i-r_j}{\norm{r_i-r_j}^3}.
\]
The symmetric initial positions are
\[
  r_1(0)=(1,0),
  \qquad
  r_2(0)=\left(-\frac12,\frac{\sqrt3}{2}\right),
  \qquad
  r_3(0)=\left(-\frac12,-\frac{\sqrt3}{2}\right).
\]

\begin{figure}[htbp]
\centering
\begin{tikzpicture}[scale=1.15, >=Latex]
  \draw[->] (-2.0,0) -- (2.2,0) node[right] {\(x\)};
  \draw[->] (0,-1.7) -- (0,1.9) node[above] {\(y\)};
  \coordinate (r1) at (1.25,0);
  \coordinate (r2) at (-0.625,1.083);
  \coordinate (r3) at (-0.625,-1.083);
  \draw[gray!70, thick] (r1) -- (r2) -- (r3) -- cycle;
  \fill[blue!70!black] (r1) circle (2pt) node[right] {\(r_1\)};
  \fill[blue!70!black] (r2) circle (2pt) node[above left] {\(r_2\)};
  \fill[blue!70!black] (r3) circle (2pt) node[below left] {\(r_3\)};
  \draw[->, red!75!black, thick] (r1) -- ++(0,0.65);
  \draw[->, red!75!black, thick] (r2) -- ++(-0.56,-0.32);
  \draw[->, red!75!black, thick] (r3) -- ++(0.56,-0.32);
  \draw[dashed] (0,0) circle (1.25);
  \node[red!75!black, align=center] at (1.45,1.2)
    {tangent velocities\\for rotating triangle};
  \node[align=center] at (0,-1.95)
    {the perturbation lab changes the initial velocities\\and monitors
     energy, angular momentum, and separation};
\end{tikzpicture}
\caption{Equal-mass three-body laboratory. The unperturbed comparison solution
starts from an equilateral triangle with tangential velocities; perturbing this
initial data is a controlled way to watch the loss of symmetric motion.}
\label{fig:lab-three-body-initial-data}
\end{figure}
\FloatBarrier

For the rotating equilateral solution, velocities are tangent to the
circumcircle. A representative perturbed initial condition uses a speed
parameter \(v=0.8\) and perturbation \(\epsilon=0.05\), with initial velocities
such as
\[
  \dot x_1(0)=\epsilon,\qquad
  \dot x_2(0)=-\frac{\sqrt3}{2}v-\epsilon,\qquad
  \dot x_3(0)=\frac{\sqrt3}{2}v.
\]
The important lesson is not the particular number \(0.8\). It is the workflow:
derive Newton's equations from a Lagrangian, integrate the exact pairwise
forces, monitor energy and angular momentum, and compare unperturbed symmetric
motion with perturbed motion. This is the finite-mass cousin of the
Sun-Jupiter-asteroid ejection simulation.

\section{Helicopter Rotor Model As A Gauge Connection}
\label{sec:lab-helicopter-gauge}

A toy helicopter model makes the Shapere-Wilczek connection finite-dimensional.
It is the most compact laboratory version of the inertia-operator connection
derived in Section~\ref{subsec:explicit-gauge-equation}: the general shape
coordinate \(s\) is replaced by two rotor angles \((\alpha,\beta)\), and the
zero-momentum equation can be solved by inspection.
Let the body have isotropic base inertia
\[
  I_{\mathrm{body}}=I_0\mathbf 1.
\]
Let one internal rotor spin about the body \(z\)-axis with angle \(\alpha\) and
inertia \(I_1\), and let a second rotor spin about the body \(y\)-axis with
angle \(\beta\) and inertia \(I_2\). The internal shape space is the torus
\[
  \mathcal S=S^1_\alpha\times S^1_\beta.
\]
Let \(\Omega\in\R^3\) be the body angular velocity. In the body frame, the
rotor angular velocities are
\[
  \Omega+\dot\alpha e_z,
  \qquad
  \Omega+\dot\beta e_y.
\]
Keeping only the rotational kinetic energy relevant for reconstruction,
\[
  T=
  \frac12I_0\norm{\Omega}^2
  +
  \frac12I_1(\Omega_z+\dot\alpha)^2
  +
  \frac12I_2(\Omega_y+\dot\beta)^2.
\]
The body angular momentum conjugate to \(\Omega\) is
\[
  \Pi
  =
  I_0\Omega
  +
  I_1(\Omega_z+\dot\alpha)e_z
  +
  I_2(\Omega_y+\dot\beta)e_y.
\]
Set total angular momentum to zero. Componentwise,
\[
  (I_0+I_1)\Omega_z+I_1\dot\alpha=0,
\]
\[
  (I_0+I_2)\Omega_y+I_2\dot\beta=0,
\]
\[
  I_0\Omega_x=0.
\]
Therefore
\[
  \Omega
  =
  -\frac{I_1}{I_0+I_1}\dot\alpha\,e_z
  -
  \frac{I_2}{I_0+I_2}\dot\beta\,e_y.
\]
In connection notation,
\[
  g^{-1}\dot g
  =
  -\mathcal A_\alpha\dot\alpha-\mathcal A_\beta\dot\beta,
\]
with
\[
  \mathcal A_\alpha=
  \frac{I_1}{I_0+I_1}e_z,
  \qquad
  \mathcal A_\beta=
  \frac{I_2}{I_0+I_2}e_y.
\]
The components are constant, but the structure group is nonabelian:
\[
  [e_z,e_y]\neq0.
\]
Thus the curvature contains the commutator term
\[
  \mathcal F_{\alpha\beta}
  =
  [\mathcal A_\alpha,\mathcal A_\beta].
\]
A small rectangular cycle in \((\alpha,\beta)\) can therefore produce a net
body rotation even though each rotor angle returns to its original value. This
is the finite-dimensional version of falling-cat holonomy.

\begin{figure}[htbp]
\centering
\begin{tikzpicture}[scale=1.0, >=Latex]
  \begin{scope}
    \draw[->] (-0.2,0) -- (4.4,0) node[right] {\(\alpha\)};
    \draw[->] (0,-0.2) -- (0,3.1) node[above] {\(\beta\)};
    \fill[blue!8] (0.85,0.65) rectangle (3.25,2.35);
    \draw[blue!70!black, thick, ->] (0.85,0.65) -- (3.25,0.65)
      node[midway, below] {\(\Delta\alpha\)};
    \draw[blue!70!black, thick, ->] (3.25,0.65) -- (3.25,2.35)
      node[midway, right] {\(\Delta\beta\)};
    \draw[blue!70!black, thick, ->] (3.25,2.35) -- (0.85,2.35);
    \draw[blue!70!black, thick, ->] (0.85,2.35) -- (0.85,0.65);
    \node[align=center] at (2.05,3.35) {closed shape loop\\on the rotor torus};
  \end{scope}
  \begin{scope}[xshift=5.4cm]
    \draw[->, thick] (0,0) -- (1.6,0) node[right] {\(e_x\)};
    \draw[->, thick] (0,0) -- (0,1.6) node[above] {\(e_z\)};
    \draw[->, thick] (0,0) -- (-1.1,-0.8) node[left] {\(e_y\)};
    \draw[red!75!black, thick, ->] (0.7,0.55) arc (30:290:0.55);
    \node[red!75!black, align=center] at (0.85,-1.2)
      {net body rotation\\from \([\mathcal A_\alpha,\mathcal A_\beta]\)};
  \end{scope}
\end{tikzpicture}
\caption{Helicopter-rotor gauge laboratory. The connection components are
constant, but they point in noncommuting body directions. A closed rectangle in
shape space therefore has curvature from the commutator term and can produce a
net body rotation.}
\label{fig:lab-helicopter-gauge-loop}
\end{figure}
\FloatBarrier

\section{What These Laboratories Practice}
\label{sec:worked-lab-summary}

\begin{itemize}
\item Coordinate changes preserve the Euler-Lagrange form when the
transformation is regular; the same physical path can have many coordinate
descriptions.
\item Regularity of the velocity Hessian is what permits a Lagrangian system to
be rewritten as a Hamiltonian system without losing information.
\item Constraints and contact forces should be solved together. Eliminating a
constraint too early can hide the force that decides loss of contact.
\item Effective potentials turn circular-orbit and bifurcation questions into
ordinary one-variable stability tests.
\item The Henon-Heiles, three-body, and rotor examples show how quickly
integrable intuition must be supplemented by phase-space geometry and
numerical diagnostics.
\end{itemize}

\section{Problem-To-Demo Conversion Rules}
\label{sec:lab-demo-rules}

Every computational laboratory in this repository should record the following
metadata:
\[
  \text{model}+\text{variables}+\text{parameters}
  +\text{units}+\text{invariants}+\text{numerical method}.
\]
For Hamiltonian demos, the minimum checks are:
\[
  \Delta H(t),\qquad
  \Delta J(t) \text{ if angular momentum is conserved},\qquad
  \text{step-size sensitivity}.
\]
For probabilistic demos such as asteroid ejection, the output must also state:
\[
  \text{sampling distribution},\quad
  \text{random seed},\quad
  \text{time horizon},\quad
  \text{ejection criterion}.
\]
Without these, a numerical picture is not a reproducible mechanics experiment.

The same rule applies in reverse when reading a demo. The code should reveal
the equations it integrates, the units in which parameters are measured, the
events it declares, and the invariants it monitors. A plot is the last line of
the argument, not the first.
